\subsection{天使大人的家庭情况}

吃饭的时候真昼回来了，三人一起吃了晚饭。真昼虽然还是显得有些僵硬，但像刚和慧见面时那种冰冷刺骨的表情已经消失，取而代之的是有些尴尬的表情。慧似乎也是一样，理所当然地，他吃得很尴尬且拘束，反而让人觉得有些对不住他。\\

话虽如此，真昼那冰冻的气氛已经消失了，周在餐桌下偶尔握一下她的手，她似乎就能平静下来，于是没发生什么特别的麻烦就吃完了饭。\\

接着，周迅速收拾完饭后的餐具，在已经提前等待着的真昼身旁坐下。\\

「……那么，真昼似乎也冷静下来了，能让我们听听你的事了吗？如果真昼觉得难受我会喊停的，可以吗？」

「好的」\\

本来这不该由周来主导，但如果交给这两个人，估计会因为顾忌而迟迟开不了口，所以周才主动搭话。

看到慧和真昼都没有异议地点了点头，周松了一口气，再次抬起头，看向隔着矮桌对面坐在沙发上的慧。\\

「呃，从哪里开始说比较好呢」

「首先，我们想听的是，你来找真昼是为了问什么事吧。真昼也是这么想的吧？」

「是的。事情发生得太突然我感到很困惑，而且我想象不到你想问我什么、想知道些什么」\\

直到现在，真昼都不知道慧的存在。虽然她听说过也许有情人的孩子，但那也是不确定的信息。仰慕小夜为母亲的孩子，竟然会来拜访自己，这绝对是她始料未及的。

在这样的情况下，对方说有事想问自己，她会警戒也是理所当然的。

说到底，怎么看都像是慧更了解小夜的事情，为什么他还要来拜访真昼呢？\\

「……说得也是。谢谢你愿意接待我唐突的拜访」\\

慧明明知道自己对真昼来说是不速之客却没有退缩，向着周和真昼恭敬地低下了头。

如果他态度傲慢，或者带有轻蔑的意味，真昼一开始大概就不会理会他了。正因为慧虽然着急却待人真诚，真昼才会缓和态度的吧。\\

周心想，这和自己仅见过一面的小夜的形象完全联想不到一起啊，同时等待着慧继续说下去。\\

「在进入正题之前，可以让我先说明一下我和姐姐的母亲之间的关系吗？」

「……好的」\\

虽然已经有了心理准备，但周还是能看到真昼再次听到母亲这个词时，身体不禁僵硬了一下。\\

「我跟……那个，我平时称呼为母亲的那个人，并没有血缘关系。所以，本来叫你姐姐是一件很奇怪的事」

「没有血缘关系……」

「我的，那个，遗传学上的母亲，已经过世了」

「简而言之，就是坐在这里的真昼的母亲在代替你母亲的职责，是这样吗？」

「是的。为了方便我还是像以前一样称她为母亲……在我懂事的时候，母亲就已经在我的记忆里了，从那时起她就一直都在。父亲告诉我生母因病去世了，所以我知道现在的母亲是继母」\\

刚才慧说过他十三岁了。一般认为人懂事大约是在两岁半到三岁左右，所以真昼大概在七岁前后。据说小雪照顾真昼是在她很小的时候，所以时间也对得上。\\

只是，听了慧的话，更加凸显了小夜对真昼毫不在意的事实，真昼也察觉到了这一点，表情上落下了一层阴影。\\

「我是被母亲抚养长大的。在知道她是继母的情况下，过着普通的……虽然工作看起来很忙，但和周围双职工家庭的生活没什么不同，我觉得我们就是个不用怀疑什么的普通家庭」\\

那份普通，真昼究竟有多么渴望呢？周不禁感到一阵心酸，但既然真昼没有大声发作，周也不该插嘴，而且这也不是慧的错，所以周把心中的苦闷留在了心里。

慧看着真昼，露出了抱歉的神色，但从他的眼神中可以看出来，他并没有停止诉说的打算。\\

「只是，那个……我发现她在户籍上也不是我的母亲」

「也就是说并没有结婚吗？」

「是的。虽然自称是母亲，但实际上算是毫无关系的陌生人」\\

这并不是什么奇怪的事。

就真昼和朝阳的话来看，并没有听说两人离了婚。小夜只是在代替母亲的职责而已，实际上并没有和慧的父亲建立婚姻关系，就是这么回事吧。\\

「被我发现的时候，呃，她说是内室。父亲说，希望在户籍上我依然是生母的孩子。关于这点我也有我的想法，但我能看到父亲很重视死去的生母，现在的母亲也很尊重这一点，就理所当然地接受了。于是我被说服了」

「……你的父亲和那个人，是在双方同意的情况下建立关系的啊」

「至少在我看来是这样的。母亲说她是生母的熟人，还是重要的朋友。有好几张两位母亲和父亲年轻时的照片，他们看起来关系很好，这点大概不是谎言。父亲也说她们两人关系很亲密。我觉得在这件事上没有说谎的必要」

「正因为那个人和你血缘上的双亲关系亲密，所以那个人也抚养了你，是这样吗？」

「如果只是到这里为止，我也觉得没什么奇怪的。在照顾亲密朋友遗留下来的孩子的过程中……电视剧里不也有这种情节吗？……但是，前几天我发现了一件奇怪的事」

「奇怪的事是指？」

「我知道了你的存在」\\

慧静静地看着真昼。\\

「契机是我在母亲房间里找到的东西」

「找到的？」

「母亲的房间里有很多装满书和文件的文件夹。那个，因为我知道母亲这段时间都会很晚回来，所以为了找学习需要的书就没打招呼进去了，这是我的不对……但从我在找书时不小心掉落的文件夹里，掉出了照片。是姐姐你，小时候的照片」\\

或许是没想到小夜会有自己的照片，真昼的表情中渐渐渗出了惊愕之色。\\

「我对这张照片没有印象。但是，上面是个长相有些眼熟的女孩。我不由自主地盯着看，拿起了装照片的文件夹。……里面还附着一封信，或者说是一份报告书之类的东西，上面写着对她对待亲生女儿态度的指责，我这才第一次知道，母亲有一个有血缘关系的女儿」\\

真昼轻轻发出啊的声音，一定是推测出了写那封信的人，就是她至今依然仰慕着的小雪吧。

虽然是工作，但小雪从小就把真昼抚养长大，真昼也知道小雪是打心底里重视她的。工作关系想必需要汇报成果，或许在那时候她就对小夜指出了关于真昼的一些问题吧。\\

「对我来说，这简直就是晴天霹雳。我一直以为我们是非常普通的家庭，听到母亲有孩子，或者说，母亲有一个和我们不同的家庭，我感到非常震惊。我以为她只是为了重视生母而没有入籍，是个普通的未婚女性。结果，事实却是这样」\\

慧那颤抖的声音，充满了对小夜的疑虑。

那苦涩的声音中，透露出不想相信的感情，但也包含了淡淡的死心。\\

「我无法忍受，追问了回来的母亲。问她这到底是怎么回事」

「……问了之后呢？」

「她不肯回答我，只说这和我没关系，小孩子不懂。问了父亲也是一样。两个人固执地都不肯回答。……于是我产生了不信任感」\\

正如他所说的，慧的表情和声音，都毫无保留地向周他们传达着满满的疑虑。\\

「因为事实就是如此吧？明明像母亲一样对待我，但其实自己有家庭，然而却对自己的孩子一直置之不理，这种事」\\

重新被慧指出的、从别人眼中看到的现实，对真昼来说是非常残酷的。\\

「虽然不知道能不能说是普遍现象，但我觉得比起别人的孩子，优先对待有血缘关系的孩子，是理所当然的事情。除非有什么极特殊的理由，否则无视自己的孩子，不是很奇怪吗？作为一个人，作为一个母亲」\\

这是明确的放弃抚养。

虽然听说没有施加暴力，但当时真昼依然品尝了难以用笔墨形容的痛苦，如果没有亲切对待她的小雪，真昼大概已经崩溃了吧。周觉得，现在的真昼能坐在这里本身就可以说是一种奇迹，这是一个如此机能不全的家庭。\\

「对我来说，母亲是值得尊敬的存在，我也有自觉被倾注了珍贵的爱意。正因为如此，我才无法相信，因为她对待亲生女儿的做法，让我怀疑起了母亲至今为止对我的爱。然后，母亲究竟在想什么，才会做出这种事？现在，那个女儿过得怎么样……会在意这些，也是顺理成章的」

「……小夜阿姨她，没有回答你对吧？」

「嗯。所以，我无论如何都想知道母亲的意图，我又偷偷找了母亲的房间，从而知道了你的事情。那个，你可能会感到不愉快，比如你是个什么样的人，住在哪里之类的。不知道算不算幸运，不知道她是不是在在意你的事，有一份像是定期报告书之类的东西，上面有现住址和照片，所以我才能找到这里」\\

只见真昼听到「定期报告书」这个词皱起了眉头。

周之前已经从慧那里听说了进行了背景调查的事，但在真昼听来应该是第一次听到吧。

周看到真昼缩起肩膀，像是警戒一样在矮桌下握紧了拳头，动作微小到只有坐在她旁边的周能察觉。

慧似乎把她的这个样子仅仅当成了不快感，满脸歉意地垂下了眼睛。\\

「因为我擅自的想法，突然出现让姐姐感到不快，真的，非常对不起。但是，我无法压抑自己。我无法继续装作不知道，无法视而不见」\\

慧深深地低下了头，连发旋都看得一清二楚，但他似乎并不后悔，而是用平静的目光注视着真昼。\\

「我不明白。母亲为什么不去抚养你，而是抚养我。因为，这不是很奇怪吗？明明有亲生女儿，却如此珍惜像我这样没有血缘关系的孩子。这只会让人觉得，是因为有什么目的才这么做的吧？」

「我想我也会产生怀疑的，如果在你的立场上」

「在我眼里的母亲，是个合理且讨厌无用功的人。她不是那种毫无意义、随随便便就会行动的人。既然那样的母亲采取了行动，就一定有什么理由」

「……你认为不单单是出于纯粹的爱吧？」

「母亲虽然不会蔑视感情，但也不是只凭感情行动的人。我认为她对我怀有爱意和她抱有某种目的，这两者并不冲突」\\

如果能断言仅仅是单纯的爱，对慧来说也许是最好的，但从慧的立场来看，抛弃亲生孩子而去疼爱没有血缘关系的孩子，如果有正常的感性，肯定会感到恐惧的。在这里怀疑除了爱以外的目的是很正常的。\\

「能想到的可能原因就是钱了，但母亲至少有自己的公司，我知道她是个相当有钱的人。我父亲也有不少钱，说到底两人本来就是对钱不太执着的人」

「……这一点我同意。那个人，虽然在金钱计算上很精明，但对于照顾我的人，她也是毫不吝啬给予报酬的。我自己到现在为止也没有为钱困扰过，如果是为了钱，应该有更多可以削减开支的地方吧」\\

别人家里的财产状况确实不太方便详细打听，周也不打算过问，不过以前真昼曾说过，父母虽然完全不照顾她，但钱倒是给得足足的。大学的费用以及好长一段时间的生活费都完全不用愁。\\

如果真的需要钱，肯定会从没有感情的真昼的费用中削减吧。而事实上并没有那样做，反而听说还有所增加，看来在费用方面确实毫不含糊。\\

真昼说出了让慧无法安心的保证，深深地叹了一口气。\\

「所以，如果有目的的话，我觉得应该是除此之外的动机」

「是的。可是，就算我再怎么烦恼为什么，我也想不明白。我能了解到的事情有限，如果我又在探查的事被发现了，我想这次她就会彻底隐藏起来。能没被隐藏起来简直是奇迹，多亏了母亲很忙」

「这对那个人来说真是罕见的百密一疏呢」

「她为了工作和亲属关系四处奔波，看起来很忙……对我来说，这是个好机会。总之，我想趁着现在还没受到限制的时候采取行动」\\

第一次见面的时候，他说过只有今天有机会。

虽然不知道他家离这里有多远，但至少小夜不可能特意住在可能会遇到真昼的地方，而慧的父亲在了解情况后依然把小夜留在身边，看来两家是有一定距离的。\\

父母双双不在家，在朋友的帮助下并且是不被怀疑的情况，以及慧有自由时间的同时真昼也在家的情况，这种条件可不多见。正如他本人所说，现在这个时候就是绝佳的机会。\\

「关于母亲，我尊敬她，作为母亲也很喜欢她。至少，对我来说，母亲是一位好母亲。……但是，我搞不懂了。那个对我既严厉又温柔的母亲，真的是母亲真实的面貌吗？我看到的母亲，难道不是个幻影吗？至今为止我看到的母亲的身影全都是虚伪的，所谓的爱，难道从一开始就不存在吗？一想到这里，我就突然觉得好害怕」\\

他原本不知道的，是自己母亲冷酷的一面。

在窥见那一角之后，慧可能就急剧地察觉到了自己周围的环境是多么地危险。就像曾经包围真昼的环境扭曲并崩坏一样，他担心总有一天也会在自己不知道的地方崩坏。\\

「所以，我来见作为母亲女儿的你。我想问问你，对你来说母亲是一个什么样的存在，还有如果你知道的话，我想问问她为什么要抚养我」\\

这是漫长而迫切的真心话吐露。

为了解决这个对幼小身躯来说过于沉重痛苦的烦恼，慧才会来到这里。既然父母不肯告诉他，想要去问另一个当事人也是极其理所当然的。即使无法解决，他也期待着至少能抓住解开烦恼的线索，于是向真昼伸出了手。

只不过，从周看来，他的这个行动恐怕也会白费吧，这是可以预见的。

因为真昼被一切所疏远地待在这里，对家族之间的牵绊一无所知。\\

「我明白你想说的了。但是，我无法回答你」

「怎么会……」

「虽然我很想回答你，但遗憾的是我知道的事情比你想象的要少得多。……我不知道那个人为什么要抚养你，但至少在我的认知里，她就是一直沉浸在情人家里不回来」

「情——」

「不忠、出轨，这么评价比较好吗。我的亲生父亲即使知道这件事也没有采取任何行动，所以我一直以为她是公然在外面寻求爱情」\\

虽然不知道事实究竟如何，但至少在真昼的视角里，小夜在外面建立家庭，而朝阳对此默许。这是破裂的夫妻关系和亲子关系。

仿佛在呼应越来越冷的声线，真昼的表情也渐渐失去了温度。周像是要拴住她一般握住她的手，或许是心理作用，一股冰冷的触感传了过来，这更进一步地将真昼的痛苦传递给了周。\\

「……我一直以为，她是在外面生了孩子。所以我曾怀疑你是不是有一半血缘关系的弟弟。不过看了一眼发现你不像母亲，我也觉得这个可能性很低」\\

看着瞳孔微暗倒映着慧的真昼，周也再次将慧的脸放在视线的中央端详。

慧这个还带着强烈稚气的少年，至少长得不像真昼。

虽然容易被他看似聪慧的气质和眼神所隐藏，但仔细看的话，会发现他有一副有些温和柔弱的容貌。虽然只看了一眼小夜的样子，记忆有些模糊，但他的长相和小夜也不同。恐怕是像父亲或者过世的生母吧。\\

面对周和真昼的视线，慧十分尴尬地游移着目光，但他轻轻叹了口气后，稳稳地接受了这两人的视线。\\

「……在姐姐看来，母亲是那种会找情人的类型啊」

「虽然我对那个人是否有那样的感情和欲望抱有疑问，但至少听说过她在外面有伴侣的事，所以我觉得大概就是这么回事吧」

「伴侣……我想她确实是伴侣。但是，在我看来，父亲和母亲并不像是那种……那个，用爱情结合在一起的关系。该说是商业伙伴吗？虽然能看出来他们很重视对方，但看起来不像恋爱之类的东西」

「也许只是在孩子面前不表现出那种气氛而已吧」

「……也许是这样，但在我看来并非如此。这是我的感想」

「是吗」

「是的」\\

面对真昼毫无温度且冷淡的回答，慧虽然有些畏怯，但并没有退缩。这说明他就是这样一直看着自己的父母的吧。这说明，他就是这样一直和父母一起度过的。\\

「无论事实如何，对我来说都是毫无关系的事」

「毫无、关系」

「因为事实就是没关系。无论那个人和你的父亲达成了什么样的协议，我都不打算做任何事。当然，也不打算去质问。做那种事对我来说非但没有任何益处，反而只会遭受损失。对现在的我来说，只希望那个人不要干涉我，不要和我扯上关系。不管有什么理由，那些人抛弃了我，这是事实」

「抛弃了，是说」

「不仅是家务、甚至连教育都全部丢给别人，自己去组建其他家庭，这不就是抛弃了吗？现实就摆在面前，那就是无论我怎么努力，那个人的关心都只向着你，只放在了你身上」\\

真昼用比平时更加辛辣的态度淡淡诉说，那副样子与其说像受了重伤，不如说是像在抚摸虽然已经愈合但仍感到紧绷的伤痕，显得很烦躁。

在慧的心中还是现在进行时，但在真昼的心中大概已经是过去的事了。总是不停受伤，好不容易伤口愈合了，但疼痛却依然残留。

「她在金钱上让我不愁吃穿，给我安排了保姆，这两点我本身是很感谢的，但仅此而已。我无法将那个人视为母亲，也不打算那么做。那个人应该也对此心知肚明吧」

对真昼来说，代替父母的只有小雪。珍视、慈爱她的，也只有小雪。

一切都已经事到如今了……真昼的眼神传达出了这个意思。\\

「所以，事到如今就算你提起那个人的事，我也不打算怎么做。这是你家庭的问题。我认为你应该和本人谈谈，决定今后的态度」

「不好意思，我也赞成。……希望你记住这并不是你的错，但这次的事情终究是你和你母亲之间的问题吧？」

「那是……」

「从经过来看，是因为你对母亲产生了不信任感，所以来问可能知道内情的人，是这样对吧？但真昼什么都不知道，对此也无能为力。就算你再向真昼渴求什么，我想真昼也给不出答案的」\\

慧的目的是，如果真昼知道内情的话就从她那里问出来。

然而，真昼对于小夜和慧父亲之间的协议，或是隐瞒的事情一无所知。说到底，她明确知道慧的存在还是在今天，根本不可能了解那些事情。

想要从真昼那里问出对慧有益的信息是不可能的，自己不知道的事情，根本就无从告知。

更重要的是，周不想让真昼再为此伤心了。\\

「你该做的事，是和你母亲面对面谈谈吧。这是产生在你们之间的问题，应该由慧和小夜女士去解决。……而且，如果可以的话，希望你不要把真昼卷入那个问题里。我知道这不是局外人该说的话，但真昼给不出你想知道的答案。真昼不知情，要说的话，反而是真昼来质问你为什么都有可能哦」\\

到底是抱着怎样的想法，才会抛弃亲生女儿，对别人家的儿子扮演母亲的角色呢。

那个儿子都找上门来了，真昼反过来去质问的未来也是有可能发生的。只是现实中，她做出了因为不想扯上关系所以不知道也罢的选择。\\

「……是我太唐突了，对不起」

「不，我这边才是态度太不成熟了，对不起。……我不希望你误会，我对你并没有怨恨或者愤怒之类的感情。哪怕我对那个人心怀愤怒，你和那个人也是不同的人」\\

正如她本人所说的那样，真昼对慧并没有什么厌恶感。

确实，在真昼看来，慧在某种意义上处于加害者的一方，但他本人什么都不知道，也是一个被大人的事情所摆布的、立场相同的受害者。

正因为在理性上理解了这一点，真昼才没有发怒或是怨恨，只是把他当作一个突然带来麻烦的人来对待吧。\\

「……姐姐对母亲的事情」

「我已经不对那个人抱有期待了，并且很看不起她的为人。尽可能地，我不想和她扯上关系」\\

毫无商量余地般冷淡，且没有一丝迟疑的回答，大概是出自真昼的真心吧。

她已经不再寻求作为母亲的爱了。\\

「你被好好地爱着呢」

「……虽然我现在开始怀疑那是不是真的了，但我确实被珍爱着抚养长大」

「是吗。这是好事」

毫无嫉妒，无比平坦的声音。

「对你来说，她是一位好母亲呢」

「……是的」\\

看着似乎很难回答但还是给予了肯定的慧，真昼深重地叹了一口气。\\

「那种类似母性的东西，那个人原来是有的啊」\\

那声音和刚才不同，仿佛是把对慧的一点点羡慕、对小夜的失望，以及曾经自己的绝望混合在一起，稀释后涂抹在表面上一般，回响着悲凉。\\

「对不起」

「这不是你该道歉的事吧。你也是被大人们随心所欲的安排所摆布的受害者」

「不，我因为自己擅自的苦衷而伤害了你，这是事实。你可以不原谅我。请不要原谅我。因为，是我夺走了你的幸福」\\

慧似乎看穿了真昼微微动摇的感情，果断地说出了责任。

本来没有责任的他，却直截了当地对真昼说出伤害了她也是事实的话语，这副模样虽然稚嫩，却很了不起。\\

「……这么年幼的孩子带着觉悟来探求真相，那个人却认为只要盖住掩饰过去就行了呢」

「虽然被称为年幼的孩子心情有点复杂，但母亲是个一旦决定就十分固执的人。我不认为我能改变她」

「想必是呢」\\

真昼一副很了解的模样，一定是从小就感受到了小夜的固执吧。无论真昼怎样伸出手，她非但没有握住，反而将手甩开了。\\

「……可以换我问几个问题吗？」

「只要我可以回答的话」

「我有件事情想要确认。我明白那个人隐瞒了关于我的事。从你的语气来看，你的父亲也知道我的存在，并且是默许的，这样说对吗？」

「虽然我没有准确地问过，但从反应来看，我想父亲也是知情并接受了这一点的」\\

虽然早有预料，但再次从慧的口中听到，连作为局外人的周都感到胃里一阵沉重的不快。

真昼大概也没见过的慧的父亲。明知道别人的女儿会不幸，却依然优先考虑了慧。

小夜轻视真昼，反倒还能让人理解一点点。

但是，作为外人，而且自己也是失去妻子带着遗留下来孩子的人，竟然能对别人的孩子做出像对自己的孩子一样的……不，是做出更让人痛苦的事的父母真的存在，这未免也太缺乏人情味了吧。\\

「说实话，关于这一点，我觉得他作为一个人实在有问题。那是眼睁睁地看着姐姐遭遇不幸而视而不见哦？很奇怪吧」

「幸好你还能觉得这件事很奇怪」

「……对于姐姐来说，我觉得我的父母做出的事被当成不是人也无可奈何，就算你当面痛骂他们也能得到原谅」

「这种机会我不需要就是了」\\

真昼将无论如何都不想扯上关系的立场推到了最前面，慧听了紧紧咬住了嘴唇。\\

「……也就是说，关于我的事，是两家公认的放置不管对吧。既然已经是生气也没用的事，那就算了，但我实在想象不出那个人会因为恋爱之类的情感而行动，所以我认为大概还是有什么利益或目的吧」

「就算有利益或目的，放弃抚养也是鬼畜的行径吧」

「我想对他们来说一定是有那么重要的什么东西吧。不过也许只是单纯不想要我而已」\\

自己亲口说自己是不需要的，真昼的心里是什么滋味呢。

周对让真昼说出这种话的真昼和慧的父母感到了难以言喻的不快，同时也像是在温暖真昼冰冷的手一般温柔地抚摸着。\\

「好啦，慧你就专注于你自己的事情吧。真昼的事情我会负起责任来处理的。虽然我可能不太靠谱就是了」\\

关于真昼，最坏的情况只要成年了也就不需要父母的同意了，所以只要等到那个时候入籍就行了。

他是个温柔的孩子所以也在顾虑真昼，但年纪这么小如果背负了太多的东西也会被压垮的吧。周希望他首先优先考虑自己的事情。\\

「还有，我可以问一个问题吗？慧，你之后打算怎么办？」

「打算怎么办……是指？」

「我是说，你瞒着母亲来到这里，听了真昼的话反而加深了疑问和疑虑吧。你是打算继续质问母亲，还是打算闭口不言把疑虑藏在心里？」\\

周明白他来这里的理由是为了向真昼打听内情，但既然现在这个目的没能实现，为了解决疑问，大概只能去问当事人的他的父母了吧。

只是，作为周来说，并不想轻易推荐他再次去询问。\\

「……说心里话我想再问一次母亲。心里带着已经找姐姐确认过事实的想法去问」

「为了你好我姑且说一句。我不打算阻止你，但要问的话最好仔细考虑清楚再去问」\\

周明白慧露出了诧异的表情，但他还是没有改变意见的打算，他思索着该怎么说才能让他明白，同时又不会伤害到他。\\

「那个，你母亲，只要听真昼的描述，我觉得是个相当严厉的人。从你的视角来看，也是个固执的人。而且慧，你怀疑母亲是出于某种目的才抚养你的，对吧？」

「是的」

「那么如果，仅仅是如果啊，如果没有感情而是出于某种目的抚养你，但在慧你一再的追问下改变了主意，这也是有可能的。如果是为了目的而抚养的话，那并不一定需要爱吧？当你问出隐瞒的事情时，她判断你已经长大了，因此改变态度也是有可能发生的事」

「……啊」\\

听了慧的叙述，周也想相信事实并非如此，但如果因为被追问，导致小夜和慧的父亲之间立下的约定变得毫无意义，以此为契机，她的态度发生改变也不是没有可能的。\\

周只从真昼和慧的口中听说过小夜，所以无法断言，但至少，可以直接预见的是，如果直接质问，小夜和慧双方对对方的意识都会发生不小的改变。

想要知道真相的心情周也很懂，但他觉得慧没有考虑到寻求真相之后可能会带来的变化这种影响。\\

「母亲发生改变这件事，你不害怕吗？而且因为这件事，你自己看待父母的眼光也可能会发生改变，你不害怕吗？如果要再去问的话，我觉得你最好做好心理准备」

「……你的意思是让我把过去封存吗？」

「不，我觉得你选哪边都无所谓。……你可能会觉得我很无情，但你们家庭的事，不是我们的问题。比起草率行动，我觉得你还是在深思熟虑之后再做决断比较好」\\

慧现在还是初中生，肯定还在父母的庇护之下。刺激父母可能会造成的影响，至少应该在脑海里留个底比较好。

虽然出于担心对慧说到了这个地步，但对周来说慧选哪边他都不会去阻止。慧终究是外人，慧的家庭也仅仅只是慧的家庭。做选择的是慧，周认为他应该做出不会让自己后悔的选择。\\

如果这是真昼和小夜的事情，周肯定会有更多的行动，也会努力让事情朝着对真昼有利的方向发展。

虽然说是恋人但也是外人，不过真昼是一个虽然说是外人但周打算让她不再是外人的女性。

如果有什么不便，周会把她从父母身边剥离出来自己照顾。虽然也有被她照顾的感觉，但他有和她共同生活的打算。周的父母也允许了这一点。

说到底，周对慧并没有做到那种地步的感情和精力，仅此而已的差别。\\

即便如此，周也不希望他变得不幸，所以仅仅只是忠告慧要深思熟虑后自己做出选择。\\

「我——」\\

就在慧打算说出自己的决断时，慧衣服口袋里传来了一阵轻快但被闷住的响声。那熟悉的声音，在慧慌忙拿出手机后变得更加清晰。\\

「……是爸爸。不会吧，他那边应该是半夜才对啊……」\\

看着盯着屏幕呆然嘟囔的慧，周浮现出了一个念头：简直就像是算准了这个时候一样。

说那边是半夜，也就是说他认为父亲不在日本。这也让他再次认同了慧说的父母都不在家的今天这个时候是个好机会。\\

「我们要回避一下吗？」

「……不，事已至此。能请你们听一下吗？」

「哎？」

「听了你们就能明白父亲是个怎样的人了。在我面前，他真的是个普通的父亲哦。虽然在你们看来可能会觉得他是个残暴不仁的人吧」\\

露出带着苦涩笑容的慧，在捕捉到周和真昼困惑的视线后，划过了接听通话的按钮。\\

「……喂，我是慧」

『慧吗。我想着差不多是春假了就联系你了——没搞错吧？』\\

可能是开了免提模式，连周他们也能听见的音量响了起来。

听到那低沉却温和的声音，周和真昼自然地挺直了背脊，为了不被发现自己的存在而紧紧闭上嘴巴，侧耳倾听。\\

「没有，已经放春假了哦。所以我不是说去朋友家住一晚嘛。工作是不是太忙了，连这都忘记了」

『说起来好像是有这么回事。因为你不在家我才吓了一跳』

「我也吓了一跳。你什么时候回国的」

『大概几个小时前吧。因为临时变更了日程，空出了很多天，所以就先回国了。难得回来，本来打算跟慧度过之后再回去工作的……是吗，去住朋友家了啊』\\

慧的父亲听起来有些遗憾，至少能让人切实感受到他对慧的爱。如果只顾工作的话，他就不会特意回国，甚至都不会有要一起度过这片刻休息时间的念头了吧。

正如慧所说，他真的深爱着自己的儿子。\\

『自从你上了初中我就一直担心你能不能交到朋友，现在看来你交到了可以约好去对方家留宿这么亲密的朋友，真是太好了。如果可以的话，也让爸爸直接和对方打个招呼吧……』

「爸爸要是突然过来大家会吓一跳的，别来了。爸爸你知道自己的长相吗？很有威慑力的哦？突然被你这种人登门拜访，朋友的父母和奶奶都会吓到的吧？」\\

这大概是配合他的朋友的家庭结构为前提在说话吧。

周觉得慧能如此流畅地夹杂着谎言说话，还挺有胆量的，不过如果不这么做，这边也会受到影响，所以还是很感谢他能如此巧妙地圆场。\\

『我自己有自知之明，你别再说了。……朋友的父母在吗，在的话我想打个招呼』

「现在在做饭，我觉得不行。话说要是现在打招呼我会觉得非常不自在的，别了吧，我和朋友的家人也会有顾忌的嘛」\\

慧大概是觉得要求周他们充当朋友父母不太妙，便面带微笑，同时冒着些许冷汗巧妙地把话题带向了没法打招呼的方向，这让周感到有些钦佩。\\

『这样啊，那没办法了。好不容易有几天的空闲，下次有机会再说吧』\\

也许是觉得没什么特别可疑的地方，慧的父亲毫不怀疑地带着轻笑接受了。\\

「就这么办吧。虽然我觉得不用太在意就是了」

『这种事情打招呼是很关键的。慧你也要记住』

「都被你说过多少次了，我知道了」

『那就好』

「只是打电话确认一下？你还真是爱操心啊」

『是啊。在意慧有没有给人添麻烦，作为父亲是理所当然的吧？』

「我才没那么调皮呢——」

『那可说不准。慧你有时候会发挥出莫名的行动力，做出一些出人意料的事情吧？』

「……爸爸你对我太失礼了」\\

现在正是他发挥这种行动力的时候，但在当下也不好指出来，周再次紧紧闭上了嘴唇。\\

『行，只要没惹出问题就行了。我一直提心吊胆，怕你在朋友家打扰别人弄出什么事来呢』

「太不信任我了。那我差不多该挂了」

『嗯。那么，也请代我向那边的两位问好』

「真是爱操心啊，知道了知道了。爸爸你也工作辛苦了。我明天就回去，回见」\\

就这样，慧在表现得没有让父亲产生任何异样感的情况下挂断了电话，但周却对刚才的对话感到了一丝不对劲。\\

（是我的错觉吗，这个人刚才是不是认知到我们的存在了）\\

慧曾传达过朋友家有父母和祖母三位监护人在场的意思，但慧的父亲口中说的却是『两位』。特意使用两位这个词的事实，让人无论如何都觉得不对劲。

今天慧之所以来这里，是因为只有这天可以自由行动，反过来说，这天会有所行动也是可以预料到的。\\

（难道说）\\

实在是想太多了吧，周把这个一闪而过的念头从脑海中甩掉，抬头看向打完电话后似乎是因为紧张消散而把身体靠在沙发上的慧。\\

「……我父亲就是这样的人。我想比两位想象的要更普通一些哦」

「与其说是普通，不如说我觉得是一位很为儿子着想的父亲呢」

「我认为他是一位连亡母的份一起好好疼爱我的父亲。真的很忙，但也真的很珍惜我。虽然因为工作的关系经常不在家，但休息日都会像这样回家看我学习，或者带我出去玩。我认为他是一位好父亲」\\

慧打从心底里敬慕父亲，从现在的声音和表情就能看出来。周如果没有任何信息，光听刚才的电话，肯定也会毫无保留地称赞他是一位好父亲吧。\\

「正因如此，我也无法相信他竟然会对姐姐的待遇视而不见」\\

这句补充的话语，瞬间将人拉回了现实。

无论对儿子有多么温柔、多么理想的父母，对真昼来说，也不过是间接伤害自己的人而已。人有多个侧面，看到的侧面不同印象就会改变，这是常有的话题。\\

第一次知道父亲另一面的慧感到困惑也是理所当然的，这也是驱使他行动的原因吧。\\

目前看来，慧对父亲的感情更多是困惑和失望，还没到轻蔑和厌恶的程度，但这样回去之后，日后再次质问小夜的时候，事情会如何发展呢。\\

「……我说啊，慧君」\\

这时，周注意到了一个问题。\\

「在」

「这下，慧你是不是回不了家了？听那气氛，你爸爸好像在家里啊」

「啊」

「……啊——」\\

在刚才的电话里，慧的父亲是因为暂时回国发现慧不在家这个理由才联系他的。

也就是说，他在家里。

慧之所以撒谎说是去朋友家留宿，是为了让不在家显得不自然，但实际上他原本打算谈完话就回家的吧。因为周一指出来，他脸上瞬间浮现出「糟了」的焦急神情，所以这个猜测应该没错。\\

足足沉默了五秒钟的慧嘟囔着「计划被打乱了一，并伴随着一声长长的叹息，垂头丧气地耷拉下了肩膀。\\

「……对不起，我太大意了。虽然事前确认过两人都不在家，但我没想到工作提前结束他就回来了……」

「听起来像是计划有变，这也是慧你没办法解决的问题啊。……只要那个计划的变动不是因为你这次拜访造成的就好」

「你说了什么吗？」

「没，这只是我的杞人忧天而已」\\

周甚至在想，说不定对方是故意放任慧行动，然后在适当的时机打了个探听情况的电话，但特意把这种不确定的因素当成他不安的材料也不好，所以周只是把它作为众多推测之一咽进了肚子里。\\

「露宿街头就不谈了，未成年在没有监护人的情况下订酒店之类的也不行。网吧未成年人也不能进。如果没有那通电话，回去还能用朋友有事之类的借口蒙混过去呢」

「真的是太出乎意料了……怎么办啊」

「要是这样的话，哎，看来只能在我家住一晚了」\\

吃完晚饭一边聊天，时间已经到了晚上八点半，外面理所当然已经是一片漆黑。就算外面有可以住的地方，也不是小孩子能随便走动的环境，很有可能会遇到可疑人员或是被警察抓走。

在这样的情况下把他赶出去，良心上实在过不去，看来只能放弃，让他在这里住一晚了。\\

「……可以吗？」

「这样做对我的心理健康也比较好？要是就这样让你回去，出了什么事我也会难受。虽然没法怎么好好招待你，如果你不介意的话」

「从头到尾都给您添麻烦了……」

「真昼没问题吗？」

「没问题是指？」

「哎，当着本人的面这么说也不太好，不过我自认为我的家对真昼来说是一个重要的归宿」\\

虽然现在她对慧的抗拒感已经不像刚见面时那么强了，但也不可能对慧毫无芥蒂。

对真昼来说，周的家就像是她自己的家一样，是归宿之一。把外人带进那里，她会不会觉得有些不快呢？这就是周的担忧。不过这似乎是杞人忧天了，真昼让亚麻色的头发微微飘动，轻轻摇了摇头。\\

「如果屋主允许的话我觉得没问题哦。说到底，把在这个时间没法回家的孩子赶出去，反而会让我的心更痛的」\\

真昼断然表示，反倒是什么办法都不想就这样让他回去更让人讨厌。周向真昼点了点头，再次看向慧，只见慧露出了一丝松了一口气的表情，恭敬地低下头说道「虽然很抱歉，但就拜托您收留我一晚了」。
